\section{Cosmic Cycles and Ages}

The world-ages, the eternal return, the golden age and its long decline. Every old tradition that felt time as a turning wheel and every one that felt it as a marching line was reaching for something real in the way the model moves. The surprise is that both were right, and for one reason. The model has cycles at one level and forbids them at another, and the same single fact gives it both at once.

\begin{principle}
The reversible interior wants to loop. The wake marches. The thin thread of docks born at peace is the whole difference between eternal return and linear history.
\end{principle}

This section states the one fact that decides everything, ranks the seven readings of cosmic cycles by how the model carries each, and works the mechanics on the committed $\{3,4,3,4\}$ bulk. The closed-region recurrence and the wake are committed parts of the theory. The age-readings are tiered by how much extra each one leans on.

Two terms used below should be fixed first. The \emph{knit} is the model's one state-law, collide then stream. Each beat the tones inside a dock interact by a fixed reversible charge-conserving table (collide), then each tone moves one dock along its site (stream). The \emph{wake} is the model's one space-law, the growing edge where new docks are born at peace at the open frontier, every beat, forever. The knit changes what is there. The wake changes how much there is.

\subsection{The one fact that decides everything}

Two of the model's commitments pull in opposite directions, and the tension between them is the answer, not a flaw to be fixed.

The first commitment is that the bulk knit is a reversible permutation. A closed, finite, reversible system is a permutation of its own configurations, and every permutation is built out of cycles. So a closed finite region, left to itself, returns. It comes back to its starting configuration after its period and then repeats, forever. This is Poincare recurrence. At the closed-and-finite level the model is intrinsically cyclic. Eternal return is not a mystical add-on here. It is the bare mathematics of a finite reversible knit.

The second commitment is that the mesh grows, forever. New docks wake into existence at peace at the frontier every beat. The wake is committed and endless, eternal expansion, no contraction, no crunch. A system that is always getting bigger never returns to a smaller past state. The wake breaks the cycle.

\begin{theorem}[Cycle below, arrow above]
A finite region of the mesh treated as closed evolves by a reversible permutation of its tone-configurations, so it is exactly periodic and returns to any past state after its period. The open whole, fed by the wake, has a strictly enlarging configuration space, so it never returns. The model is cyclic at the closed-finite level and one-way at the open-whole level, and the wake is the sole difference.
\end{theorem}

So the precise answer to whether the model has cosmic cycles is a clean split. Yes for any closed finite region, an exact cycle, Poincare recurrence. No for the open growing whole, where the wake turns the cycle into a one-way history.

\subsection{Time, precisely}

Time in the model is the wake. It is the advance of the open edge, the radial direction along which the mesh gets larger. The present is the wake. Time is not a dimension the world sits inside. It is the world getting bigger.

This is exactly why the whole cannot cycle. A cycle of the whole would require the world to get smaller again, to un-grow, to un-wake docks back into nothing. The model forbids that. The wake only ever runs one way.

\begin{principle}
Cyclic time and linear time are the same time seen at two levels. The reversible interior wants to loop. The wake marches.
\end{principle}

\subsection{The readings of cosmic cycles, ranked}

Every traditional cycle-idea is one of the seven readings below. They are ordered by how cleanly the committed model carries them, and each carries a status tier.

\subsubsection{Closed-region recurrence, the exact eternal return}

Any finite region treated as closed, with no wake crossing its boundary, is a reversible permutation, so it cycles exactly. It returns to its initial state after its period, then repeats forever. This is the eternal return, made exact and provable.

The catch is the period. It is the cycle length of a permutation on a huge state space, so it is astronomically long. The return is real in principle and never seen in practice. And it holds only for closed regions. The open growing world does not have it.

\begin{result}[Status, Tier 1]
Closed-region recurrence is exact and grounded. It is the bare mathematics of a reversible finite knit. The tradition it fits is Nietzsche's eternal recurrence and the cyclic cosmos, true locally, false for the growing whole.
\end{result}

\subsubsection{The dynamic-balance cycle, the ages as order and erosion}

The lean creates order out of peace, and the churn of the reversible interior erodes it, and a region swings between the two. This is the dynamic balance the arrow holds. A region rises to high order, a golden age. It then erodes through declining order, the ages falling away. It bottoms out in disorder, a dark age. Then it is regenerated from peace by the lean, and the turn begins again.

The decline of the four classical ages, with their lengths in the ratio four, three, two, one, is the erosion phase. Order wears down at an accelerating relative pace. The turn of the age is the lean rebuilding from the still point. This is the cleanest reading of the four ages specifically. They are the phases of one order-and-erosion cycle.

\begin{result}[Status, Tier 1]
The dynamic-balance cycle is grounded in the lean and the dynamic balance the arrow holds. The tradition it fits is the ages of gold down to iron, the decline and renewal of the world's quality.
\end{result}

\subsubsection{Cycles within cycles, the nested timescales}

The nesting of an age within a greater age within a still greater one is the integration tower, the nested selves. Each level has its own lifecycle at its own timescale. A self's life sits inside a culture's life inside a world's life, each a cycle of order and erosion at its own pace, the smaller turning many times within one turn of the larger. The cosmic cycles-within-cycles are simply the multi-scale temporal structure of the integration hierarchy.

\begin{result}[Status, Tier 1]
Nested cycles are grounded in the multi-level selves. The tradition it fits is the nested ages and epochs, the wheels within wheels.
\end{result}

\subsubsection{Scale recurrence, the cycle as a holographic layer}

The hyperbolic bulk is self-similar across scale. Each radial level is a scaled copy of the last, the mipmap of the holographic renormalization. So a kind of cycle is the scale-step. It is not a return in time but a recurrence in scale, the same structure appearing again one level out, larger.

On this reading the ages are not later, they are bigger. Each great age is a scale-octave of the one before. This is a real recurrence, but in the radial-scale direction, not in time.

\begin{result}[Status, Tier 1]
Scale recurrence is grounded in the holographic self-similarity of the bulk. The tradition it fits is the worlds proceeding from worlds, the fractal cosmos, the as-above-so-below.
\end{result}

\subsubsection{Local dissolution and recreation, breathing without a crunch}

Worlds and selves form, when a pattern integrates, and dissolve, when it loses integration back into the churn, over and over, while the mesh itself persists and grows. The breathing of the cosmos, the day and night of creation, is read as the local formation and dissolution of integrated structures, not a global contraction. The night is the unintegrated churn between integrated days. This keeps the breathing real without breaking eternal expansion.

\begin{result}[Status, Tier 2]
Local dissolution and recreation is grounded in self-formation and dissolution, but it is gated by persistence, so it sits one tier down. The tradition it fits is creation and dissolution, the breathing of the cosmos.
\end{result}

\subsubsection{The double-cover cycle, the two-turn return}

The bulk has a literal built-in two-cycle. A spinor returns to itself only after two full turns, through $720$ degrees rather than $360$, the spin-half double cover. So there is an exact cyclic structure in the geometry where one round is not a return and two rounds are. It is a faint, precise echo of the idea that a single age does not complete the pattern and a pair does, the world returning after the second turn, not the first.

\begin{result}[Status, exact geometry]
The double-cover cycle is exact geometry, a curiosity more than a cosmology. The tradition it fits is the doubled ages, the idea that completion needs a second pass.
\end{result}

\subsubsection{Limit-cycle worlds}

The coarse-grained dynamics of a region could settle onto a periodic attractor, a limit cycle, so a world runs through a fixed periodic sequence of states. This is possible but not forced. It depends on the region's effective dynamics, and the open-system growth of the wake tends to perturb exact periodicity.

\begin{result}[Status, Tier 3]
Limit-cycle worlds are plausible but dynamics-dependent. The tradition it fits is the wheel of fixed recurring ages.
\end{result}

\subsubsection{What the model rules out}

Three traditional readings are forbidden outright, each by the wake.

\begin{table}[ht]
\centering
\begin{tabular}{@{}ll@{}}
\toprule
reading & why the model forbids it \\
\midrule
a global crunch and rebirth & eternal expansion, the wake never contracts \\
exact recurrence of the whole universe & the world is always larger than its past, it cannot return \\
time as a closed loop with no progress & the wake is a real arrow, history does not exactly repeat \\
\bottomrule
\end{tabular}
\caption{The cycle-readings the model rules out, and the reason in each case.}
\label{tab:ruled-out}
\end{table}

\subsection{How it works on the $\{3,4,3,4\}$ bulk}

The readings run with real mechanics on the committed substrate, not as analogies but as statements about the actual lattice.

A finite $\{3,4,3,4\}$ torus, the $24$ sites with the reversible pair knit and no wake, is a permutation of its tone-configurations, so it has exact cycles. Run it long enough and it returns. This is the closed-region recurrence, reading one, exact on the actual substrate.

Add the wake, new docks born at peace at the frontier, and the configuration space keeps enlarging, so the global return is destroyed and an arrow appears. The wake is the exact mechanism that turns the cyclic torus into an arrowed history.

The order-and-erosion cycle, reading two, runs as the dynamic balance on the bulk. The lean seeds order at the peace wake, the reversible interior churns it, and the coarse-grained order of a region rises and falls.

The scale recurrence, reading four, is the radial mipmap of the bulk, exact, each Busemann level a scaled copy. The nesting, reading three, is the coarse-graining tower over the bulk, each level a slower clock.

\begin{proposition}[What the substrate carries]
On the committed $\{3,4,3,4\}$ substrate the model carries readings one through five with real mechanics, reading six as exact geometry, and reading seven as a possibility. It rules out the global crunch and the exact global return.
\end{proposition}

\subsection{The wheel with a leak}

The deepest reading is that the model holds cyclic time and linear time together, and they are the two faces of its one mechanism. The reversible bulk is a cycle. The wake is an arrow. The thin thread of docks born at peace is the whole difference.

\begin{principle}
Time is a cycle that never quite closes, because the world is always a little bigger than it was.
\end{principle}

The traditions that felt time as a wheel were glimpsing the reversible cyclic interior. The traditions that felt time as a line, a history, a progress, were glimpsing the wake. Both are true, from their angles, and the model says exactly why. The wheel is the closed mathematics. The line is the open wake. The real world is the wheel with a leak.

\begin{principle}
The age turns, the world returns to a golden age, but never to the same golden age, because the crystal that turns has grown in the turning.
\end{principle}

\subsection{The honest bottom line}

Cosmic cycles are real at the closed-and-finite level, exact Poincare recurrence, provable, Tier 1. They are real as an order-and-erosion rhythm, the lean's dynamic balance, Tier 1. They are real as nested timescales and as scale self-similarity, Tier 1. The local breathing rests on persistence and so sits at Tier 2. The double cover is an exact geometric curiosity. The limit cycle is dynamics-dependent and sits at Tier 3.

Forbidden is the global return and the contracting crunch, ruled out by eternal expansion.

\begin{result}[The cleanest single statement]
The world is a reversible cycle with an unstoppable wake, so it forever tries to return and forever fails to, and that failing-by-a-little is the arrow of time. The ages are real, the wheel is real, and the wheel never closes, because the world keeps growing through it.
\end{result}
